\subsection{Adversarial evidence and the physical-truth boundary}
A receipt can bind an authorized recorder, a mission, source, a plan and an outcome digest. Verification checks those predicates. It cannot establish that a physical sensor produced the claimed observation. The signed-registry controls reject modified payloads, cross-mission replay, expiry, a wrong purpose, an unknown signer, revocation and a changed historical ledger entry. A deliberately false claim signed by the trusted local recorder passes signature verification. This is an expected and necessary negative control.

The finite completion inspector checks the supplied model's allocation and objective. The independent outcome evaluator checks synthetic truth. In the drift condition, these two checks can disagree about usefulness without a software-integrity defect. An attacker who can supply a coherent false world and regenerate its proofs has crossed the study's trusted-input boundary. We do not claim trusted hardware, remote attestation, an independently calibrated sensor, or an external witness.

\subsection{The no-NFT counterfactual}
The local-EVM demonstration binds a compact summary of the actual completed development pilot: six worlds and 42 method records, including the exact raw-data digest. It is a separate mechanism control, not an on-chain execution of the 3,072-world confirmation. The off-chain registry binds the same evidence bytes, plan and base-seed parent commitment and checks once-only routing. Both preserve the same productive result because neither changes the instrument allocator or its physical-model execution. This equality is structurally expected: it tests replacement of the studied evidence-binding interface, not market coordination or every institutional benefit of a public registry.

The local lifecycle executes 113 transactions, including deployment, fixture setup and negative cases. Total simulated gas is 20,701,896; the separately tracked lifecycle subset uses 4,546,309 gas. These are local EVM quantities, not a gas-price quote or a live financial bill. The off-chain record occupies 3,619 bytes. Its construction and repeated verification timings are retained with the source-bound run. These resources cannot be subtracted from one another as though they were a common cost unit.

An off-chain signature requires an independently trusted public key and retained state. The demonstration rejects tampering and an omitted middle record. It also accepts a valid truncated prefix when the verifier has no separately pinned terminal checkpoint. A locally operated EVM can likewise be reset or presented as a different history; this experiment does not confer public-network consensus. Actual public anchoring could provide shared ordering and enforce escrow under additional assumptions. The off-chain control does not replicate decentralized escrow, token custody or public finality.

\scope{\textbf{Registry result.} The NFT layer contributes no measured instrument-output gain in the functional counterfactual. The signed registry is a credible simpler design for the studied evidence-binding service. This does not establish equivalence for every contract function or prove that on-chain accountability has no institutional value.}

\subsection{Open knowledge and limited disclosure}
The open branch publishes algorithms, source, reproduction instructions and synthetic inputs. The restricted branch uses Node/OpenSSL AES-256-GCM with a fresh ephemeral 256-bit key, a fresh 96-bit nonce and mission-associated authenticated data \citep{dworkin2007}. Both voluntarily authorized operators receive the same plaintext; outsider access is restricted. Authorized decryption, wrong-key rejection, changed ciphertext rejection and changed-mission rejection are tested. No encryption key is persisted. Metadata and length remain visible. Since the generator and final synthetic workloads are public, this demonstrates disclosure-control functionality, not secrecy of the published experiment.

A small-domain commitment test recovers an unsalted SHA-256 calibration label by enumerating 16 possibilities. Hashes are commitments, not encryption. Withholding a key is an access restriction and is never counted as superior productive capability. Ed25519 is the existing signature primitive, not new cryptography \citep{rfc8032}.

\subsection{Five contracts, verified source semantics}
All five requested contracts are byte-identical between \href{https://github.com/MontrealAI/AGI-Alpha-Agent-v0/tree/197147ffd22c5ae13113162afffaf3ce9c10aa2e/contracts/ascension}{observed Alpha source \code{197147ff...}} and the full historical reference \code{c0fc5217e364a39a4cd123d17a8317c06b80f820}. This is a source comparison, not a fresh verification of live bytecode or activation. The complete source hashes, exact line links, compiler policy and local outputs appear in the companion audit.

\textbf{NovaSeed.} The ERC-721 stores capsule, plan and evidence commitments, a creator, a parent and a job count. Commitments must be nonzero; the job bound is 1--128; the metadata URI has a bounded length. A nonzero parent must be owned by the caller. A commitment does not encrypt its contents, validate a plan's productive merit, guarantee availability, or confer intellectual-property rights.

\textbf{AscensionAccess and AscensionRiskOracle.} Access checks current ENS ownership and bounded governance admission. Risk review uses eligible staked validators, a configured quorum and time-limited approval. Parent approval is not successor approval. Address separation is not proof of beneficial or institutional independence. Admission restrictions are permissions, not measurements of ability.

\textbf{AscensionMark.} An eligible seed owner opens a funding campaign with bounded target and timing. Funding lots provide internal accounting; they are not equity or a dividend promise. Funding-lot purchases and transfers have phase/time restrictions. Unused lot sales retain a withdrawal path while the phase is Funding, including after its funding deadline. Bloom requires the Funding phase, sufficient reserve, current owner/admission/risk conditions and a timestamp strictly before \code{fundingEnd}. It changes phase to Bloomed, sets the execution end and emits an event. Bloom neither burns the NFT nor mints a new agent, generates profit or locks general token circulation.

The Merkle plan commits the job index and specification. Routing checks membership, bounds and the unused index; the once-only flag is set before the external job post. The stored JobMarket specification hash uses index zero, while the Mark route event preserves the indexed plan leaf. These are distinct hashes for later indices. Unused or returned reserves can be reclaimed according to closure and pro-rata accounting; executed expenditure does not guarantee principal recovery.

\textbf{AscensionJobMarket.} The code implements opted-in staked roles, per-bid collateral, three selected reviewers and exact submitted-artifact hashes. Two approving votes can settle the job; failure, refund and collateral consequences follow their specified paths. The gross review pool is 5\%, and each actual payout burns its integer-floor 1\% share. That payout mechanism is separate from Bloom. Local execution exercises these semantics with mock assets, never live funding or trade. The Manager USDC route remains a different denomination and accounting path.
